% ===========================================================================
% CHAPTER 7
% ===========================================================================

\chapter{Certificate Analytics and Dashboard}
\label{chap:analytics-api-dashboard}

\section{Analytics Processing Framework}
\label{sec:generic-precomputation-framework}

The analytical layer combines materialized result documents with live queries
over the certificate dataset. Materialized processing creates the indexes
required by the analytical workloads and then computes six scope-aware result
families. Live queries support certificate lists, detailed records, selected
counts, risk candidates, and time-series views that depend on the current
dataset.

Figure~\ref{fig:analytics-materialization-flow} illustrates the relationship
between the two query paths.

\begin{figure}[H]
  \centering
  \begin{tikzpicture}[
    font=\scriptsize,
    >=Latex,
    box/.style={draw, rounded corners=2pt, align=center, text width=2.55cm, minimum height=0.8cm, fill=black!5},
    job/.style={draw, rounded corners=2pt, align=center, text width=1.55cm, minimum height=0.75cm, fill=black!3},
    store/.style={draw, cylinder, shape border rotate=90, aspect=0.22, align=center, text width=3.2cm, minimum height=1.0cm, fill=black!8},
    flow/.style={-{Latex[length=2mm]}, line width=0.5pt},
    live/.style={-{Latex[length=2mm]}, line width=0.5pt, dashed}
  ]
    \node[store] (source) at (0,0) {Main certificate\\dataset};
    \node[box] (indexes) at (0,-1.35) {Create and verify\\analytical indexes};

    \node[job] (ca) at (-5.0,-2.8) {CA};
    \node[job] (geo) at (-3.0,-2.8) {Geography};
    \node[job] (san) at (-1.0,-2.8) {SAN};
    \node[job] (keys) at (1.0,-2.8) {Shared\\keys};
    \node[job] (sig) at (3.0,-2.8) {Signature/\\hash};
    \node[job] (valid) at (5.0,-2.8) {Validity};

    \node[store] (results) at (0,-4.35) {Materialized global and\\country-level results};
    \node[box] (api) at (0,-5.85) {Application queries\\and filtering};
    \node[box] (ui) at (0,-7.2) {Interactive dashboard\\and drill-down};

    \draw[flow] (source) -- (indexes);
    \foreach \n in {ca,geo,san,keys,sig,valid} {
      \draw[flow] (indexes) -- (\n);
      \draw[flow] (\n) -- (results);
    }
    \draw[flow] (results) -- (api);
    \draw[flow] (api) -- (ui);
    \draw[live] (source.east) -- ++(2.3,0) |- node[pos=0.72,right,align=left]{live records, counts,\\risk, and trends} (api.east);
  \end{tikzpicture}
  \caption{Materialized and live analytics paths}
  \label{fig:analytics-materialization-flow}
\end{figure}

Materialized results represent the latest completed analytical refresh, while
live queries can observe the current main dataset independently. All analytical
families support a global view and configured country-level domain scopes.
Table~\ref{tab:analytics-family-contracts} summarizes the principal
capabilities.

\begin{table}[H]
  \centering
  \small
  \setlength{\tabcolsep}{4pt}
  \caption{Certificate analytics capabilities}
  \label{tab:analytics-family-contracts}
  \begin{tabularx}{\textwidth}{@{}>{\RaggedRight\arraybackslash}p{0.20\textwidth}Y>{\RaggedRight\arraybackslash}p{0.25\textwidth}@{}}
    \toprule
    \textbf{Family} & \textbf{Principal analysis} & \textbf{Delivery mode} \\
    \midrule
    Validity & Lifetime statistics, expiry status, duration buckets, issuance, and near-expiry windows & Materialized summaries with selected live counts \\
    Geography & Distribution derived from domain suffix and configured scope & Materialized by scope \\
    Signature and key & Signature algorithms, hash families, key algorithms, key sizes, and self-signed records & Materialized summaries with live encryption counts \\
    SAN & SAN totals, count groups, wildcard use, suffix groups, and certificate references & Materialized summaries and drill-down \\
    Shared keys & Distinct certificates grouped by repeated public-key fingerprint & Materialized groups, metadata, and detail views \\
    Certificate Authority & Issuer distributions, validation levels, self-signed counts, and comparative ranking & Materialized by scope \\
    Trends and risk & Issuance, expiration, algorithm, validation, key-size, and security-indicator views & Live or combined queries \\
    \bottomrule
  \end{tabularx}
\end{table}

\section{Certificate and Security Analytics}
\label{sec:validity-geographic-analytics}

\textbf{Validity and geographic analysis.}
Validity processing converts the encoded certificate lifetime into days and
calculates average, minimum, maximum, and duration-group statistics. It also
supports issuance and expiration timelines together with live 7-, 30-, and
90-day expiry windows. Geographic categories are derived from the domain
suffix through a configured mapping. They describe domain namespace and must
not be interpreted as verified physical location.

\textbf{Signature, hash, and encryption analysis.}
Certificates are grouped by signature algorithm, hash family, key algorithm,
and key size. SHA-1 is treated as deprecated, while MD5 and MD2 are treated as
critical hash choices. Live encryption views count selected RSA and ECDSA
algorithm-size combinations. The resulting score and strength labels are
project-defined analytical indicators rather than standardized security
ratings.

\textbf{SAN analysis.}
The SAN family calculates total and average SAN counts, single-name and
multi-name distributions, wildcard occurrence, count buckets, and common
suffix groups. Bounded certificate references connect aggregate groups to
paginated certificate drill-down.

\textbf{Shared-key analysis.}
Shared-key detection first collapses repeated observations of the same
certificate fingerprint and then retains a public-key fingerprint only when it
is associated with more than one distinct certificate fingerprint. Each group
records affected certificates, domains, SANs, issuers, key characteristics,
and summary counts. Exposure labels use the number of certificates and SANs to
support prioritization, but a shared key does not by itself establish
private-key compromise or unsafe cross-owner use.

\textbf{Trend analysis.}
Live aggregations present recent issuance activity, upcoming encoded expiry
dates, signature-algorithm adoption, validation levels, and selected key-size
series. Expiration views group scheduled \texttt{notAfter} values; they are not
predictive models.

\section{Certificate Authority Analytics and Ranking}
\label{sec:ca-analytics-ranking}

Certificate Authority analysis groups leaf certificates by issuer
organization, validation level, self-signed status, and issuer-related
characteristics. Comparative ranking is based on five component groups shown
in Table~\ref{tab:ca-score-components}.

\begin{table}[H]
  \centering
  \small
  \setlength{\tabcolsep}{4pt}
  \caption{Certificate Authority ranking components}
  \label{tab:ca-score-components}
  \begin{tabularx}{\textwidth}{@{}>{\RaggedRight\arraybackslash}p{0.25\textwidth}>{\RaggedRight\arraybackslash}p{0.29\textwidth}Y@{}}
    \toprule
    \textbf{Component} & \textbf{Principal inputs} & \textbf{Purpose} \\
    \midrule
    Core hygiene & ZLint errors, warnings, and selected critical findings & Measures certificate-quality findings \\
    Cryptographic health & Key characteristics, algorithm hygiene, and public-key reuse & Represents cryptographic and key-management indicators \\
    Operational stability & Certificate authority and operational indicator values & Captures operational characteristics defined by the model \\
    Policy compliance & Extended key usage, policy identifiers, validation level, and name constraints & Represents policy-related certificate attributes \\
    Risk factors & Issuer-country category, authority-access information, and revocation-service presence & Adds contextual risk and revocation indicators \\
    \bottomrule
  \end{tabularx}
\end{table}

For certificate $i$, the ranking score is the mean of the five normalized
component values:

\begin{equation}
  S_i = 100 \times
  \frac{H_{\mathrm{core}} + H_{\mathrm{crypto}} + H_{\mathrm{operational}}
  + H_{\mathrm{policy}} + H_{\mathrm{risk}}}{5}.
  \label{eq:ca-score}
\end{equation}

The score for a Certificate Authority is obtained by averaging the scores of
its included leaf certificates. The output also records the number of
certificates contributing to the ranking and the component averages. The
formula supports transparent comparison within the platform but has not been
externally calibrated as an industry security rating.

\section{Security-Oriented Risk Assessment}
\label{sec:vulnerability-risk-model}

The risk view forms a bounded union of certificates matching security-relevant
signals, removes duplicate records, applies the score in
Table~\ref{tab:vulnerability-score}, and ranks the resulting candidates.

\begin{table}[H]
  \centering
  \small
  \setlength{\tabcolsep}{4pt}
  \caption{Application-defined certificate risk score}
  \label{tab:vulnerability-score}
  \begin{tabularx}{\textwidth}{@{}>{\RaggedRight\arraybackslash}p{0.31\textwidth}>{\Centering\arraybackslash}p{0.15\textwidth}Y@{}}
    \toprule
    \textbf{Signal} & \textbf{Points} & \textbf{Condition} \\
    \midrule
    Expired certificate & $+30$ & Certificate status is expired \\
    Shared public key & $+30$ & Public-key fingerprint belongs to a shared-key group \\
    Weak encryption & $+20$ & RSA modulus length is below 2,048 bits \\
    Long validity & $+10$ & Encoded certificate lifetime exceeds 398 days \\
    Certificate-quality findings & $+1$ to $+10$ & Errors and warnings contribute up to ten points \\
    Positive characteristic & $-5$ each & Current validity, strong key, shorter lifetime, and absence of lint findings \\
    \bottomrule
  \end{tabularx}
\end{table}

The final value is clamped to the interval $[0,100]$ and classified as
\texttt{Critical}, \texttt{High}, \texttt{Medium}, or \texttt{Low}. Because
candidate queries are bounded and the weights are defined by the application,
the result is a prioritization aid rather than proof of vulnerability,
compromise, or exhaustive risk coverage.

\section{API and Interactive Dashboard}
\label{sec:django-api-design-filtering}

The application service exposes structured access to dashboard summaries,
scope selection, certificate lists and details, filters, analytical families,
shared-key groups, risk candidates, and trends. Certificate-list filtering
supports combinations of status, issuer, country scope, text search,
cryptographic attributes, validity windows, SAN characteristics, shared-key
membership, and risk signals.

The dashboard uses a typed client and shared state to coordinate scope,
filters, search, pagination, theme, loading state, and drill-down navigation.
Figure~\ref{fig:frontend-component-flow} summarizes the presentation
architecture.

\begin{figure}[H]
  \centering
  \begin{tikzpicture}[
    font=\scriptsize,
    >=Latex,
    box/.style={draw, rounded corners=2pt, align=center, text width=2.6cm, minimum height=0.9cm, fill=black!5},
    wide/.style={draw, rounded corners=2pt, align=center, text width=4.0cm, minimum height=0.9cm, fill=black!3},
    flow/.style={-{Latex[length=2mm]}, line width=0.5pt}
  ]
    \node[wide] (root) at (0,0) {Application layout and shared providers};
    \node[box] (scope) at (-4.2,-1.6) {Scope, search, filters,\\and page state};
    \node[box] (pages) at (0,-1.6) {Analytical pages\\and detail views};
    \node[box] (storage) at (4.2,-1.6) {Local and session\\state persistence};
    \node[wide] (components) at (0,-3.2) {Cards, charts, tables, pagination, and certificate drill-down};
    \node[box] (client) at (-3.0,-4.8) {Typed request client\\and page adapters};
    \node[box] (cache) at (0,-4.8) {Client request reuse\\and revalidation};
    \node[box] (service) at (3.0,-4.8) {Application service\\with scoped queries};

    \draw[flow] (root) -- (scope);
    \draw[flow] (root) -- (pages);
    \draw[flow] (root) -- (storage);
    \draw[flow] (scope) -- (components);
    \draw[flow] (pages) -- (components);
    \draw[flow] (storage) -- (components);
    \draw[flow] (components) -- (client);
    \draw[flow] (components) -- (cache);
    \draw[flow] (client) -- (service);
    \draw[flow] (cache) -- (service);
  \end{tikzpicture}
  \caption{Interactive dashboard component relationships}
  \label{fig:frontend-component-flow}
\end{figure}

Major views cover platform health, active and expired certificates, validity,
signature and hash distributions, CA analysis and ranking, SAN analysis,
shared-key groups, trends, country-level exploration, security-oriented risk,
and detailed certificate inspection. Aggregate charts and tables link to
filtered certificate lists or individual records, providing a consistent path
from high-level observations to supporting certificate data.
